\chapter{Четырёхвершинное вектороподобное замыкание и дефицит селектора}
\label{ch:v7-four-vertex-vectorlike-selector-gate}

\section{Кандидат после аномального запрета}

Минимальный консервативный ремонт добавляет полный синглетный и дублетный
вектороподобные пакеты
\begin{equation}
 X_L,X_R\sim(\mathbf1,\mathbf1)_{-1},\qquad
 Y_L,Y_R\sim(\mathbf1,\mathbf2)_{-1/2}.
 \label{eq:v7-vectorlike-four-vertex-content}
\end{equation}
Левые и правые вклады сокращаются попарно, поэтому
\begin{equation}
 \mathcal A_{221}=\mathcal A_{\mathrm{grav}^2 1}
 =\mathcal A_{111}=0,
 \qquad N_{\mathbf2}^{\rm new}=0\pmod2.
 \label{eq:v7-vectorlike-four-vertex-anomalies}
\end{equation}
Вектороподобные расширения конечных спектральных троек действительно могут
иметь калибровочно-инвариантные массы и быть безаномальными
\cite{v7vector-classification,v7vector-scalars}. Но это ещё не выбирает их
связи с исходным поколением.

\section{Полный первопорядковый граф}

Для девяти физических вершин полный перебор критерия одной строки или одного
столбца даёт
\begin{equation}
 |E_{\rm strict}|=14,\qquad
 |E_{\rm old}|=3,\qquad
 |E_{\rm new}|=11.
 \label{eq:v7-vectorlike-edge-count}
\end{equation}
Желаемый шестикромочный цикл вместе с двумя собственными
вектороподобными массами использует только шесть новых блоков:
\begin{equation}
 |E_{\rm selected}\setminus E_{\rm old}|=6,
 \qquad
 |E_{\rm allowed}\setminus E_{\rm selected}|=5.
 \label{eq:v7-vectorlike-edge-deficit}
\end{equation}
Оставшиеся разрешённые рёбра включают дополнительный кварк-лептонный канал,
смешивание старых и новых копий и ещё один внутренний путь. Первый порядок
их разрешает, но не заставляет зануляться.

При трёх поколениях каждое новое ребро несёт независимую комплексную матрицу
$3\times3$. До факторизации базисных вращений это означает
\begin{equation}
 11\times 9\ \text{комплексных}
 =198\ \text{вещественных параметров}.
 \label{eq:v7-vectorlike-raw-yukawa-cost}
\end{equation}
Следовательно, аномальный ремонт сам по себе возвращает произвольное меню
юкавских матриц, против которого был сформулирован стоп-критерий предыдущей
развилки.

\section{Почему лёгкий хиральный сектор всё же не исчезает}

В заряженном синглетном секторе общая семейная масса имеет размер
$3\times6$, а в дублетном --- $6\times3$:
\begin{equation}
 M_e:\mathbb C^6_R\to\mathbb C^3_L,\qquad
 M_L:\mathbb C^3_R\to\mathbb C^6_L.
 \label{eq:v7-vectorlike-mass-shapes}
\end{equation}
Для общей матрицы обе имеют ранг три, поэтому
\begin{equation}
 \dim\ker M_e=3,\qquad
 \dim\ker M_L^*=3.
 \label{eq:v7-vectorlike-kernel-count}
\end{equation}
То есть на каждое поколение остаются один правый заряженный синглет и один
левый слабый дублет. Их направления, однако, являются произвольными смесями
старых и новых копий.

На уровне хиральных индексов результат точен:
\begin{equation}
 I_{\rm four\ vertex}=(1,-1,-1,1,-1)=I_{H_{15}}
 \label{eq:v7-vectorlike-index-preservation}
\end{equation}
в порядке $(Q,u,d,L,e)$. Полное зеркальное поколение вместо этого даёт
\begin{equation}
 I_{\rm full\ mirror}=(0,0,0,0,0),
 \label{eq:v7-vectorlike-full-mirror-index-zero}
\end{equation}
и при общей массе уничтожает весь хиральный индекс. Поэтому полный зеркальный
атом является более крупным, но не более подходящим ремонтом.

\section{Дубликаты бимодулей и отсутствующий центр}

Расширенный носитель содержит совпадающие типы
\begin{equation}
 L_L\simeq Y_L=(\mathbb H,\mathbb C,L),\qquad
 e_R\simeq X_R=(\mathbb C,\mathbb C,R).
 \label{eq:v7-vectorlike-duplicate-types}
\end{equation}
Неизменённая алгебра действует на каждую пару одинаково и потому допускает
унитарные вращения пространства кратностей. Она не различает «старую» и
«новую» копии, не выбирает три лёгких ядровых направления и не выделяет
шесть желаемых блоков из одиннадцати разрешённых.

\begin{theorem}[Условный допуск четырёхвершинного замыкания]
Четырёхвершинный вектороподобный сектор проходит строгий первый порядок,
Real-замыкание, локальные и глобальную аномальные проверки. В отличие от
полного зеркального поколения он сохраняет хиральные индексы $H_{15}$ и
оставляет по одному лёгкому лептонному направлению на поколение. Однако
неизменённая алгебра не выбирает ориентацию этих ядер и разрешает пять
лишних рёбер. Поэтому получен допустимый носитель, но не канонический
single-source родитель.
\end{theorem}

\begin{nogocriterion}[Запрет произвольного юкавского ремонта]
Нельзя вручную занулить пять разрешённых блоков или выбрать ядра массовых
матриц по наблюдаемым смешиваниям. Требуемый селектор должен следовать из
уже существующего аффинного/Hodge-родителя либо быть объявлен новой
архитектурой до вычисления спектра.
\end{nogocriterion}

\begin{protocol}
\begin{enumerate}
 \item строгий первый порядок: \textbf{пройден};
 \item Real и аномалии: \textbf{пройдены};
 \item хиральный индекс $H_{15}$: \textbf{сохранён};
 \item полное зеркальное поколение: \textbf{отклонено индексом};
 \item новых разрешённых матричных блоков: \textbf{$11$};
 \item блоков желаемого ремонта: \textbf{$6$};
 \item лишних разрешённых блоков: \textbf{$5$};
 \item ориентация лёгкого ядра выведена: \textbf{нет};
 \item итог: \textbf{условно положительный носитель, селектор отсутствует};
 \item следующий гейт: \textbf{селектор пространства кратностей из
 аффинного/Hodge-родителя}.
\end{enumerate}
\end{protocol}

Машинный сертификат сохранён в
\path{s2t_v7_four_vertex_vectorlike_selector_gate_results.json}.

\begin{thebibliography}{9}
\bibitem{v7vector-classification}
J.-H.~Jureit, T.~Schuecker, C.~A.~Stephan,
\emph{On a Classification of Irreducible Almost-Commutative Geometries V},
Journal of Mathematical Physics 50 (2009), 072303; arXiv:0901.3214.
\bibitem{v7vector-scalars}
C.~A.~Stephan,
\emph{New Scalar Fields in Noncommutative Geometry},
Physical Review D 79 (2009), 065013; arXiv:0901.4676.
\end{thebibliography}